% Generated draft: requires mathematical and source review.
\documentclass[11pt]{article}
\usepackage[margin=25mm]{geometry}
\usepackage{fontspec}
\setmainfont{Latin Modern Roman}
\usepackage{amsmath,amssymb,mathtools,unicode-math}
\setmathfont{Latin Modern Math}
\usepackage{enumitem,needspace,xurl,hyperref}
\setlength{\parindent}{0pt}
\setlength{\parskip}{5pt}
\setlength{\emergencystretch}{3em}
\newcommand{\N}{\mathbb{N}}
\newcommand{\Z}{\mathbb{Z}}
\newcommand{\Q}{\mathbb{Q}}
\newcommand{\R}{\mathbb{R}}
\newcommand{\C}{\mathbb{C}}
\newcommand{\F}{\mathbb{F}}
\newcommand{\A}{\mathbb{A}}
\newcommand{\PP}{\mathbb{P}}
\newcommand{\E}{\mathbb{E}}
\newcommand{\eps}{\varepsilon}
\newcommand{\dd}{\,\mathrm{d}}
\begin{document}
\section*{UnsolvedMath ID: 30003106}
\textbf{Source problem code:} catalogue entry
\subsection{Definitions and mathematical statement}
For an integer \(q\geq 2\) and a coefficient field \(K\subseteq\C\), a convergent power-series germ \(g\) at \(0\) is called a \(q\)-Mahler function over \(K\) if it satisfies an identity
\[
\sum_{j=0}^{r}a_j(z)g(z^{q^j})=0,
\qquad a_j(z)\in K(z),
\]
with the coefficients not all zero. A usual system presentation for a vector \(F=(f_1,\ldots,f_m)^{\mathsf T}\) is
\[
F(z)=A(z)F(z^q),\qquad A(z)\in\operatorname{GL}_m(K(z)).
\]
The supplied record does not specify \(K\), the analytic hypotheses, or this precise system convention. These definitions provide a conditional framework, rather than additional hypotheses attributed to the source.

For this system convention, one usual meaning of a regular point \(\beta\) is that neither \(A\) nor \(A^{-1}\) has a pole at any point \(\beta^{q^k}\), for every integer \(k\geq0\). Let \(S\) be the set of algebraic points \(\beta\) satisfying \(|\beta|<1\), at which the specified function values exist and which are nonsingular in the intended sense. The record does not establish whether nonsingularity means this entire-orbit condition or only regularity at \(\beta\). The boundary \(|\beta|=1\) is excluded. The record does not exclude \(\beta=0\), so zero must remain eligible whenever the chosen hypotheses permit it.

\textbf{Requested determination, subject to these unresolved specifications.} Given the vector of \(q\)-Mahler functions and an algebraic point \(\alpha\in S\), determine, for a specified transcendental Mahler function \(g\), the set
\[
E(g;S)=\{\beta\in S:g(\beta)\in\overline{\Q}\}.
\]
Here the function \(g\), its relationship to the given vector, and the base field relative to which it is transcendental must be specified; transcendence over \(K(z)\) means that no nonzero polynomial in \(K(z)[X]\) vanishes at \(g\).

If the requested ``value relations'' mean polynomial relations with algebraic coefficients among the entries of \(F(\alpha)\), their precise object is the ideal
\[
I_{\alpha}(F)=\{P\in\overline{\Q}[X_1,\ldots,X_m]:
P(f_1(\alpha),\ldots,f_m(\alpha))=0\}.
\]
If instead the target is to locate points realizing a specified relation \(P\), the corresponding set is
\[
Z_P(F;S)=\{\beta\in S:P(f_1(\beta),\ldots,f_m(\beta))=0\}.
\]
These are explicit possible meanings of the requested generalization. The source does not select between them, specify the allowed relations, or define what constitutes a determination procedure. Thus it does not supply a unique fully specified mathematical problem beyond this qualified formulation.
\subsection{Short English statement}
Determine which algebraic points strictly inside the unit disk give algebraic values of a transcendental Mahler function, subject to the intended nonsingularity condition, and determine the associated value relations. The coefficient field, the role of the given point \(\alpha\), the meaning of nonsingularity, and the precise relation problem remain unspecified. The displayed sets and ideal express conditional interpretations requiring human mathematical review.
\subsection{Sources}
\begin{itemize}
\item Local UnsolvedMath registry, supplied complete source record: permanent ID \(30003106\), exact source problem code \texttt{catalogue entry}, title \emph{Algebraic Values and Relations of Mahler Functions}; collection: Oberwolfach Reports Open Problems. No individual proposer is supplied. The record gives proposed year \(2016\).
\item Supplied report citation: \emph{Diophantische Approximationen (2017)}, Oberwolfach Reports, DOI \texttt{10.4171/owr/2016/21}, pp.~1099--1187. Supplied source URL: \url{https://doi.org/10.4171/owr/2016/21}. This is report-level provenance; no problem-specific page or named contribution is supplied.
\item Supplied literature source, labeled ``OWR report'': \url{https://ems.press/content/serial-article-files/46625?nt=1}.
\item Literature assessment dated exactly \texttt{2026-08-21}: the supplied curation labels the problem open and states that the \(q\)-Mahler transcendence-degree theorem does not determine the requested higher-order relation points, and that no general determination procedure is established in the record or checked literature. This is supplied dated curation, not an independently verified assessment. The background explicitly requests independent verification of time-sensitive status and cited claims.
\item The registry labels the statement \texttt{corrected\_verified} and says its mathematical content was preserved while notation, typography, or omitted definitions were repaired. These are supplied editorial assertions, not mathematical verification of this draft.
\end{itemize}
\end{document}
